% !TeX root = ../main.tex
% ============================================
% 第一章 导论
% ============================================

\chapter{导论}

\section{研究背景与意义}

\subsection{主动健康管理中的连续风险评估需求}

慢性病通常具有起病隐匿、病程较长和影响因素交织等特点，已成为公共健康领域长期面对的问题\cite{gbd2019global}。以糖尿病、心血管疾病和代谢综合征为代表的常见慢病，其风险并不只取决于某一项异常指标，而往往受到遗传背景、生活方式、临床指标变化和环境暴露等因素共同影响。对于这类疾病，仅依赖一次体检结果或阶段性门诊记录，往往只能看到某一时点的状态，难以支撑对个体风险演化过程的持续判断。

现有医疗服务在诊断和治疗环节已经形成较成熟的工作流程，但整体上仍以症状驱动和就诊驱动为主，更适合处理已经出现明确表现的健康问题。对于尚处于风险累积、指标波动或生活方式持续偏离阶段的人群，这种模式提供的信息通常较为离散，既不容易呈现风险随时间变化的趋势，也不便于比较干预前后的变化情况。

随着可穿戴设备、家庭健康终端、基因检测以及数字病历技术逐步普及，个体健康数据的获取频率和类型都在增加\cite{sockolow2021mhealth,zahedani2023digital,tang2024ehr}。这使健康管理系统不再只是保存体检记录，而需要进一步承担多源数据组织、风险分析和后续干预支持等任务。主动健康管理关注的，正是如何利用这些持续积累的数据，对尚未发展为明确临床事件的风险状态进行更早、更连续的识别与跟踪。

\subsection{多模态融合与智能交互的现实挑战}

尽管数字健康工具不断普及，健康管理场景中仍然普遍存在数据分散、接口割裂和语义不统一的问题。体检系统通常侧重结构化临床指标，基因检测服务更多以静态报告形式输出结果，可穿戴设备主要产生心率、步数和睡眠等连续观测数据，智能问答又依赖知识检索和上下文组织来给出解释性回复。不同模态在采样频率、表示方式和可信度上的差异，使这些数据很难直接组合成统一、可计算的个体画像\cite{ambalavanan2025ecosystem,carter2022ehrgenomics,epstein2023ehrwearable}。

这一问题直接影响平台侧的分析能力。单一模态往往只能反映个体健康状态的局部信息，难以同时覆盖遗传易感性、当前表型和行为变化；即便已经采集到多源数据，如果缺少统一的接入、组织和更新机制，这些信息也不容易转化为稳定且便于解释的分析结果。对慢性病管理而言，平台如果停留在“记录”和“展示”层面，就很难支撑连续风险评估。

此外，健康管理并不止于识别风险，还要求系统能够在风险评估之后支持方案比较，并在新的观测数据到来后更新分析结果。用户在实际使用中也往往希望把档案数据、文档内容、历史结果和知识解释放到同一个交互入口中理解，而不是在多个工具之间反复切换。基于这一现实需求，本文将问题聚焦于如何在同一平台内组织多模态输入、综合风险分析、情景化干预推演和受控问答流程。

\subsection{研究意义}

本文的研究意义主要体现在工程组织和应用支撑两个层面。工程上，论文并不将多模态健康管理拆分为彼此孤立的模型或工具，而是尝试把数据接入、风险分析、干预推演和问答交互放在同一平台中讨论，使各模块之间的数据流和责任边界更清楚。对于需要同时处理结构化指标、文档内容、知识检索结果和交互状态的健康管理系统，这种组织方式有助于后续实现、联调和验证。

在应用层面，本文实现的原型为体检指标、遗传信息、文档解析结果和行为相关信息的联合使用提供了一个可运行的样例，可用于支撑风险识别、方案比较和咨询辅助等任务。需要说明的是，本文工作的目标是形成一个与现有实现一致、能够被测试和分析的工程原型，而不是对临床效果或大规模落地价值作超出当前证据范围的推广判断。

\section{国内外研究现状}

\subsection{多模态健康数据融合研究现状}

多模态健康数据融合关注的是如何把来源不同、结构不同、时间粒度也不同的健康信息组织到同一分析框架中。近年来，相关研究首先把重点放在数据互联和标准化组织上。例如，面向数字健康数据集成生态的综述工作指出，跨系统互操作、主键映射、数据治理和长期维护机制仍然是平台建设中的基础约束\cite{ambalavanan2025ecosystem}。围绕临床数据与基因数据的协同使用，也已有研究从电子病历互通、基因结果入库和语义标准对齐等角度进行了讨论\cite{carter2022ehrgenomics}。

另一条研究路径更关注不同来源数据在具体场景中的联合使用。电子病历与可穿戴数据的关联研究表明，连续行为观测能够为传统临床记录提供额外时间维度的信息\cite{epstein2023ehrwearable}。这类工作说明，多源数据并置本身具有实际价值，但也进一步暴露出一个问题，即不同模态之间的数据质量、缺失模式和解释口径并不一致。很多已有研究能够完成某两类数据的连接，却未必进一步解决统一分析入口、持续更新和结果解释如何落到同一系统中的问题。

因此，当前多模态健康数据融合研究已经从“能否接入”逐步转向“接入之后如何稳定分析”。对于本文所关注的主动健康管理场景而言，真正的难点不只是把数据收进来，而是让结构化指标、文档内容、遗传信息和行为观测能够在同一平台内进入同一条分析链路。

\subsection{疾病风险预测与评估研究现状}

疾病风险预测是健康智能分析中较早成熟的一条研究方向。对于结构化健康数据，梯度提升树模型仍然是常见技术路线。LightGBM 和 XGBoost 在表格数据建模、缺失值处理和特征贡献分析方面具有较好的工程适应性，因此被广泛用于慢性病筛查和风险分类任务\cite{ke2017lightgbm,chen2016xgboost}。在遗传信息利用方面，多基因风险评分方法为个体易感性量化提供了较稳定的补充路径\cite{choi2019prsice2}。

不过，从现有工作来看，临床风险建模和遗传风险建模常常分别展开，然后在结果层面做相对松散的解释性拼接。对于长期健康管理而言，这种做法能够回答“当前风险有多高”，却不一定能很好回答“风险由哪些来源共同影响”以及“行为调整后风险可能怎样变化”。换句话说，许多研究更重视预测性能本身，而较少继续向面向干预的比较分析延伸。

这也说明，面向平台化应用的风险评估并不只是把模型接到前端展示页面上。它还需要考虑不同证据来源如何共同进入分析过程，以及分析结果是否能够继续服务于后续的方案比较和用户解释。本文的风险分析部分正是沿着这一思路展开，即在临床指标基础上引入遗传和行为相关修正，并把结果继续连接到干预推演环节。

\subsection{智能健康管理平台研究现状}

随着移动互联网、云服务和个人健康设备的发展，健康管理平台已经从单纯记录工具逐步演化为同时承担数据采集、行为提醒、长期随访和轻量咨询的综合系统。围绕慢病自我管理的综述研究表明，数字平台在改善依从性、增加行为反馈频率和支持长期随访方面具有实际应用潜力\cite{sockolow2021mhealth}。也有研究尝试把可穿戴数据与行为模式分析结合起来，用于代谢健康管理等具体场景\cite{zahedani2023digital}。

但从平台能力形态看，现有系统往往各有侧重。一类平台长于设备接入、数据展示和趋势可视化，适合承担个人健康记录中枢的角色；另一类平台更重视咨询交互或行为建议，但对个体画像、文档内容和历史分析结果的联动支持相对有限。真正把多模态接入、风险分析、方案比较和交互解释放在同一系统里统一组织的工作还不多。

这意味着健康管理平台研究已经不再只是前端界面或单点功能的问题，而越来越接近“系统如何组织分析链路”的问题。对于本文的实现目标来说，平台层的关键不是简单叠加模块，而是让上传、解析、分析、比较和问答这些环节在数据和接口层面能够接续起来。

\subsection{医疗智能问答与Agent应用研究现状}

大语言模型的发展推动了智能问答从单轮文本生成走向检索增强和工具增强。在医学场景中，已有研究显示大语言模型能够编码一定临床知识，但其回答仍需要证据、流程和安全边界约束\cite{singhal2023clinicalllm}。ReAct、Toolformer 等工作说明，模型如果能够在回答过程中结合检索、工具调用和中间决策，往往比纯生成式问答更适合处理需要外部证据支撑的任务\cite{yao2023react,schick2023toolformer}。这一思路也逐渐影响到医疗问答系统的设计，使问答模块不再只是“回答一句话”，而是开始承担检索、判断和调用外部能力的协调角色。

不过，医疗场景对 Agent 的要求明显高于一般信息问答。系统不仅要回答得通顺，还要处理安全分流、证据充分性、工具权限边界和过程留痕等问题。已有研究已经指出，医疗多 Agent 系统在协作和安全控制方面存在额外风险，说明“能用工具”与“能在医疗场景中稳定使用”之间仍有距离\cite{chen2025medsentry}。

因此，当前医疗智能问答研究给出的启发主要在于两点：一是问答需要与检索和工具结合，二是这种结合必须受到明确约束。相比之下，如何把 Agent 与个体健康画像、历史分析结果以及平台内已有服务稳定衔接起来，并通过测试和行为评测形成可回归的运行链，仍然是值得继续完善的部分。

综合来看，现有研究已经分别在多模态数据组织、疾病风险预测、数字健康平台和工具增强问答方面积累了较多成果，但这些成果往往分布在不同层面。对于主动健康管理场景，真正困难的地方在于如何把数据接入、风险分析、情景比较和受控交互放进同一条系统链路中，并保持实现边界清晰、结果可解释、流程可验证。本文的工作正是在这一问题背景下展开。

\section{研究内容与技术路线}

\subsection{研究内容概述}

围绕前文提出的问题，本文的研究内容围绕“多模态输入组织、综合风险分析、结果延伸与反馈”这条主链展开。首先，针对健康数据来源分散、表达形式不一的问题，本文围绕结构化体检指标、遗传信息、OCR 医疗文档内容以及行为相关信息，设计统一的数据接入与组织方式，使这些异构信息能够进入同一分析流程，为后续综合风险分析提供共同的数据基础。

在核心分析层，本文以结构化健康档案为主要分析入口，在临床风险估计基础上引入遗传因素和生活方式相关修正，形成面向个人健康管理的综合风险结果。进一步地，针对用户在健康管理中常见的“如果调整当前方案，风险会怎样变化”这一需求，本文在同一风险引擎基础上引入情景化推演机制。当前实现主要围绕体重变化这一可控变量构造生活方式干预场景，用于支持未来风险模拟和方案比较。

在结果解释与系统落地层，本文结合检索增强与受控 Agent 流程，构建面向健康咨询场景的问答能力，使问答过程能够利用用户上下文、知识检索结果和只读工具输出，同时保留证据约束与安全分流逻辑。在此基础上，论文进一步完成前后端分层平台原型实现，并通过回归测试、浏览器场景验证和 Agent 行为评测，对主要功能链路进行检查和分析。

\subsection{技术路线}

本文的技术路线可概括为“多模态输入组织 $\rightarrow$ 综合风险分析 $\rightarrow$ 情景化推演 $\rightarrow$ 受控问答与结果反馈”。其中，第一步是建立统一输入层，将用户档案中的结构化指标、上传文档的 OCR 处理结果、遗传信息以及行为相关状态组织为可供后续服务调用的数据表示。第二步作为平台的分析中枢，在这一基础上执行综合风险分析，以临床风险结果为起点，结合遗传和生活方式修正，输出可解释的综合风险结果。第三步是在既有风险结果之上加入情景扰动机制，用于比较不同干预条件下的风险变化趋势，当前实现主要围绕体重变化这一生活方式干预场景展开。

在上述分析主链之后，本文将问答模块定位为分析结果的解释与反馈入口，而不是独立于平台之外的附加功能。问答流程会结合用户画像、历史分析结果、知识检索内容以及受限工具调用结果生成回答，并通过证据约束和安全分流控制输出边界。最终，平台以前后端分层实现这些能力，并通过功能测试与行为验证检查从数据输入到结果反馈的主要运行链路。

\section{主要创新点}

结合本文的实现范围，本文工作的首要特点在于，并不是围绕单一模型或单一功能展开，而是把结构化健康档案、遗传信息、OCR 文档内容和行为相关状态统一组织到同一综合风险分析入口下，使多模态输入能够在平台内进入同一条可计算主链。对于主动健康管理这类需要跨越多种数据来源和服务边界的系统，这一点构成了全文最核心的工程贡献。

在这一主链基础上，本文进一步将风险结果向两个方向延伸。一方面，平台在共享风险引擎上加入情景化推演能力，当前实现主要围绕体重变化这一生活方式干预场景，用于支持方案比较和趋势解释；另一方面，平台在分析结果之上接入受控 Agent 流程，将检索、只读工具调用、安全分流和行为验证放入同一运行约束下，使问答模块承担结果解释与反馈角色。需要说明的是，这里的“创新”主要指面向当前论文目标的系统整合、链路组织与可验证实现，而不是脱离现有代码和验证证据去宣称更强的理论突破或临床结论。

\section{论文组织结构}

全文共分为七章。第1章为导论，说明研究背景、研究现状、论文内容与技术路线，并对全文结构作总体介绍。第2章介绍论文所依赖的理论基础和关键技术，包括多模态健康数据融合、疾病风险预测与情景推演，以及 OCR、RAG 和 Agent 问答相关方法。第3章给出系统需求分析与总体设计，说明平台的功能需求、总体架构、数据模型及关键模块划分。

第4章进一步讨论平台的关键算法与模型设计，重点包括数据预处理、疾病风险预测、多模态融合风险评估以及干预推演方法。第5章对应系统实现，说明后端核心服务、前端主要功能和 Agent 相关实现。第6章通过实验与测试对平台的主要功能链路进行验证和分析。第7章对全文工作进行总结，并结合当前实现边界讨论不足与后续改进方向。
